\subsection{Complete Metric Spaces and Completions}\label{subsec:complete-metric-spaces-and-completions}

Completeness asks whether Cauchy approximations converge inside the ambient space.\\
This subsection records standard complete spaces, the contraction principle, and completion by Cauchy sequences.

\begin{example} \label{ex:finite-dimensional-complete-metric-spaces} \hlt{Finite-Dimensional Complete Metric Spaces}\\
Let $k\in\N$. For $\mathbf{z},\mathbf{w}\in\C^k$, set $d_2(\mathbf{z},\mathbf{w})\defeq\bigl(\sum\limits_{j=1}^k\abs{z_j-w_j}^2\bigr)^{1/2}$. Then $(\R^k,d_2)$, $(\C^k,d_2)$ are complete; case $k=1$ includes $\R$ and $\C$.
\begin{proof}
Space $\R^k$ is complete by Corollary \ref{cor:euclidean-cauchy-convergence}. Write $z_j=x_j+iy_j$, $w_j=u_j+iv_j$, define $F:\C^k\to\R^{2k}$ by $F\mathbf{z}\defeq(x_1,y_1,\ldots,x_k,y_k)$.\\
$d_2(F\mathbf{z},F\mathbf{w})^2=\sum\limits_{j=1}^k[(x_j-u_j)^2+(y_j-v_j)^2]=d_2(\mathbf{z},\mathbf{w})^2$ by Definition \ref{def:complex-absolute-value-analysis}.\\
The coordinate map $F$ is bijective, and the distance identity makes it an isometry by Definition \ref{def:isometry-and-isometric-metric-spaces}(i).\\
If $(\mathbf{z}_n)$ is Cauchy in $\C^k$, then $(F\mathbf{z}_n)$ is Cauchy in $\R^{2k}$ and converges to some $\mathbf{u}\in\R^{2k}$.\\
The identity $d_2(\mathbf{z}_n,F^{-1}\mathbf{u})=d_2(F\mathbf{z}_n,\mathbf{u})\to0$ implies $\mathbf{z}_n\to F^{-1}\mathbf{u}$. Hence $\C^k$ is complete by Definition \ref{def:complete-metric-space}.\\
\hlt{Structural Commentary}: The coordinate isometry transfers real finite-dimensional completeness to complex coordinates.\\
\hlt{Key Insight}: Separating real and imaginary parts preserves the Euclidean distance.
\end{proof}
\end{example}

\begin{definition} \label{def:contraction-mapping} \hlt{Contraction}\\
Let $(X,d)$ be a metric space. A map $\varphi:X\to X$ is a \hlt{contraction} iff $(\exists c\in[0,1))(\forall x,y\in X)\,d(\varphi(x),\varphi(y))\leq c\,d(x,y)$.
\end{definition}

\begin{theorem} \label{thm:contraction-principle} \hlt{Contraction Principle}\\
If $X$ is a nonempty complete metric space and $\varphi:X\to X$ is a contraction, then $(\exists!\,x\in X)\,\varphi(x)=x$.
\begin{proof}
\begin{enumerate}[label=(\roman*)]
\setlength{\itemsep}{0pt}
\item \hlt{Iteration estimate}: Choose $x_0\in X$ and a contraction constant $c\in[0,1)$ by Definition \ref{def:contraction-mapping}.\\
Set $x_{n+1}\defeq\varphi(x_n)$. The identity $d(x_1,x_0)=c^0d(x_1,x_0)$ is the induction basis.\\
If $d(x_{n+1},x_n)\leq c^nd(x_1,x_0)$, then $d(x_{n+2},x_{n+1})\leq c\,d(x_{n+1},x_n)\leq c^{n+1}d(x_1,x_0)$ by the contraction inequality.\\
Hence $(\forall n\in\N_0)\,d(x_{n+1},x_n)\leq c^nd(x_1,x_0)$ by Theorem \ref{thm:mathematical-induction} with index $n+1$.
\item \hlt{Cauchy control}: If $0\leq n<m$, then $d(x_n,x_m)\leq\sum\limits_{j=n}^{m-1}d(x_j,x_{j+1})$ by the triangle inequality.\\
By (i), $\sum\limits_{j=n}^{m-1}d(x_j,x_{j+1})\leq d(x_1,x_0)\sum\limits_{j=n}^{m-1}c^j$. The bound $d(x_n,x_m)\leq\frac{c^n}{1-c}d(x_1,x_0)$ follows by $(1-c)\sum\limits_{j=n}^{m-1}c^j=c^n-c^m\leq c^n$.\\
The right side tends to $0$ by Theorem \ref{thm:special-sequence-limits}(v) and Theorem \ref{thm:complex-sequence-limit-algebra}(ii). Hence $(x_n)$ is Cauchy by Definition \ref{def:cauchy-sequence}.
\item \hlt{Fixed-point existence}: By Definition \ref{def:complete-metric-space}, choose $x\in X$ with $x_n\to x$. Bound $d(\varphi(x_n),\varphi(x))\leq c\,d(x_n,x)$ implies $\varphi(x_n)\to\varphi(x)$.\\
The shifted sequence satisfies $x_{n+1}\to x$ by Corollary \ref{cor:metric-convergent-subsequences}. Since $\varphi(x_n)=x_{n+1}$, $\varphi(x)=x$ by Theorem \ref{thm:metric-sequence-basic-properties}(ii).
\item \hlt{Uniqueness}: If $\varphi(x)=x$ and $\varphi(y)=y$, then $d(x,y)\leq c\,d(x,y)$, hence $d(x,y)=0$ and $x=y$.
\end{enumerate}
\hlt{Structural Commentary}: Iteration turns the contraction estimate into a geometric Cauchy estimate.\\
\hlt{Key Insight}: Completeness supplies the limiting point; the strict constant $c<1$ forces uniqueness.
\end{proof}
\end{theorem}

Figure \ref{fig:contraction-fixed-point-iteration} combines contraction of distances with convergence of the orbit.
\begin{illustrationfigure}[H]
  \centering
  \begin{tikzpicture}[
  scale=0.90,
  >=Stealth,
  every node/.style={font=\small},
  frame/.style={draw=black!60, line width=0.6pt},
  image/.style={fill=gray!17, draw=black!60, line width=0.7pt},
  inward/.style={->, black!60, line width=0.65pt},
  iterate/.style={black!75, line width=0.45pt},
  point/.style={circle, fill=black, inner sep=1.25pt}
]
  \draw[frame] (-3.6,-2.08) rectangle (3.6,2.03);

  \draw[image]
    plot[smooth cycle, tension=0.72] coordinates {
      (-2.05,-1.35) (-2.48,-0.58) (-2.38,0.42) (-1.92,1.18)
      (-0.78,1.34) (0.28,1.22) (1.25,1.06) (2.28,0.82)
      (2.70,0.06) (2.55,-0.76) (1.66,-1.12) (0.62,-1.22)
      (-0.45,-1.14) (-1.35,-1.20)
    };

  \node at (2.78,1.47) {$M$};
  \node[fill=gray!17, inner sep=1pt] at (1.45,0.52) {$\varphi(M)$};

  \draw[inward] (-3.06,0.82) -- (-2.42,0.55);
  \draw[inward] (-3.05,-0.05) -- (-2.48,-0.05);
  \draw[inward] (-2.75,-1.30) -- (-2.18,-1.03);
  \draw[inward] (0.12,1.78) -- (0.00,1.28);
  \draw[inward] (3.04,0.96) -- (2.44,0.70);
  \draw[inward] (3.08,-0.24) -- (2.67,-0.24);
  \draw[inward] (1.68,-1.72) -- (1.30,-1.17);

  \node[point,label=left:$x$] at (-2.95,0.18) {};
  \node[point,label=above:$y$] at (0.70,1.48) {};
  \node[point,label=above left:$\varphi x$] at (-0.98,0.18) {};
  \node[point,label=above:$\varphi y$] at (0.58,0.54) {};

  \node[point,label=below:$x_0$] (x0) at (-1.20,-1.48) {};
  \node[point,label=below right:$x_1$] (x1) at (-0.72,-0.68) {};
  \node[point] (x2) at (-0.28,-0.50) {};
  \node[point] (x3) at (0.05,-0.38) {};
  \node[point,label=above:$x_n$] (xn) at (0.28,-0.31) {};
  \node[point,label=right:$p$] (p) at (0.70,-0.12) {};

  \draw[iterate] (x0) -- (x1) -- (x2) -- (x3) -- (xn) -- (p);
\end{tikzpicture}

  \caption{The iterates approach the fixed point $p$ while $\varphi$ contracts pairwise distances.}
  \label{fig:contraction-fixed-point-iteration}
\end{illustrationfigure}

\begin{perspective} \label{per:contractions-feed-local-analysis} \hlt{Fixed Points as Existence Proofs}\\
Theorem \ref{thm:contraction-principle} solves $x=\varphi(x)$ by convergence of iterates in a complete space.\\
This viewpoint later motivates local inverse-function arguments but does not replace their hypotheses or proof.
\end{perspective}

\begin{example} \label{ex:rational-metric-not-complete} \hlt{A Noncomplete Metric Space}\\
The metric space $(\Q,d)$ with $d(x,y)\defeq\abs{x-y}$ is not complete.
\begin{proof}
Let $(r_n)$ be decimal truncations $1.4,1.41,1.414,1.4142,\ldots$ of $\sqrt{2}$. Then $(\forall n\in\N)\,r_n\in\Q$, $r_n\to\sqrt{2}$ by Proposition \ref{prop:decimal-truncations-squeeze-real}(ii).\\
The sequence $(r_n)$ is Cauchy in $\R$ by Theorem \ref{thm:cauchy-convergence-criterion}(i), hence also in $\Q$ with the inherited metric.\\
If $(\Q,d)$ were complete, choose $r\in\Q$ with $r_n\to r$ by Definition \ref{def:complete-metric-space}. The same convergence holds in $\R$.\\
Thus $r=\sqrt{2}$ by Proposition \ref{prop:convergent-sequence-basic-consequences}(i), contrary to Proposition \ref{prop:rational-square-root-two-gap}(i).\\
\hlt{Structural Commentary}: The decimal approximants are Cauchy before their missing limit is available inside $\Q$.\\
\hlt{Key Insight}: The inherited metric preserves Cauchy control but not the irrational limit.
\end{proof}
\end{example}

\begin{perspective} \label{per:cauchy-construction-of-reals-parallel-model} \hlt{Cauchy Construction as a Parallel Model}\\
Example \ref{ex:rational-metric-not-complete} exhibits a rational Cauchy sequence with a missing endpoint.\\
For the set $\mathcal{C}_\Q$ of rational Cauchy sequences, define $(a_n)\sim(b_n)\Longleftrightarrow\abs{a_n-b_n}\to0$.\\
The quotient $\mathcal{C}_\Q/{\sim}$ motivates the construction below; only the proved completion theorem supplies later authority.
\end{perspective}

Figure \ref{fig:cauchy-completion-formal-limits} depicts the quotient relation and the embedded constant sequences.
\begin{illustrationfigure}[H]
  \centering
  \begin{tikzpicture}[
  >=Stealth,
  every node/.style={font=\small},
  stepbox/.style={rounded corners=2pt, align=center, text width=2.95cm, minimum height=1.05cm, inner xsep=6pt, inner ysep=4pt},
  box/.style={stepbox, draw=gray!60, fill=gray!7},
  classbox/.style={stepbox, draw=blue!65!black, fill=blue!7},
  realbox/.style={stepbox, draw=orange!75!black, fill=orange!10},
  arrow/.style={->, line width=0.55pt, gray!70},
  include/.style={->, dashed, line width=0.5pt},
  note/.style={font=\scriptsize, align=center, text=gray!65}
]
  \node[box] (seq) at (0,0) {Cauchy sequence\\$(a_n)$ in $\Q$};
  \node[classbox] (eq) at (3.85,0) {same real when\\$|a_n-b_n|\to0$};
  \node[realbox] (real) at (7.7,0) {real number\\$[(a_n)]=\operatorname{LIM} a_n$};

  \draw[arrow] (seq) -- (eq);
  \draw[arrow] (eq) -- (real);

  \node[box] (const) at (0,-1.65) {rational $q$\\as $(q,q,q,\ldots)$};
  \node[classbox] (constclass) at (3.85,-1.65) {constant class\\$[(q,q,\ldots)]$};
  \node[realbox] (embed) at (7.7,-1.65) {embedded rational\\$q=\operatorname{LIM} q$};

  \draw[arrow] (const) -- (constclass);
  \draw[arrow] (constclass) -- (embed);
  \draw[include, blue!65!black] (const.north) -- (seq.south);
  \draw[include, blue!65!black] (constclass.north) -- (eq.south);
  \draw[include, orange!75!black] (embed.north) -- (real.south);
\end{tikzpicture}

  \caption{Equivalent rational Cauchy sequences determine one formal limit, while constant sequences embed $\Q$.}
  \label{fig:cauchy-completion-formal-limits}
\end{illustrationfigure}

\begin{definition} \label{def:metric-completion} \hlt{Metric Completion}\\
A \hlt{completion} of a metric space $(X,d)$ is a triple $(Y,d_Y,\tau)$ satisfying:
\begin{enumerate}[label=(\roman*)]
\setlength{\itemsep}{0pt}
\item $(Y,d_Y)$ is complete;
\item $\tau:X\to Y$ is an isometry;
\item $\overline{\tau(X)}^{\,Y}=Y$.
\end{enumerate}
\end{definition}

\begin{theorem} \label{thm:metric-completion-existence-uniqueness} \hlt{Existence and Uniqueness of Metric Completions}\\
Let $(X,d)$ be a metric space.
\begin{enumerate}[label=(\roman*)]
\setlength{\itemsep}{0pt}
\item The space $(X,d)$ has a completion.
\item Let $(Y,d_Y,\tau)$ and $(Z,d_Z,\sigma)$ be completions of $(X,d)$. Then $(\exists!\,\Phi:Y\to Z)$ such that $\Phi$ is a bijective isometry and $\Phi\circ\tau=\sigma$.
\end{enumerate}
\begin{proof}
\begin{enumerate}[label=(\roman*)]
\setlength{\itemsep}{0pt}
\item \hlt{Existence}:
\begin{enumerate}[label=(\alph*)]
\setlength{\itemsep}{0pt}
\item \hlt{Pseudometric}: Let $\operatorname{CS}(X)$ be the set of Cauchy sequences in $X$.\\
For $a=(a_n),b=(b_n)\in\operatorname{CS}(X)$, set $D(a,b)\defeq\lim\limits_{n\to\infty}d(a_n,b_n)$. This limit exists by Proposition \ref{prop:cauchy-sequences-bounded-distance-converges}(ii).\\
The map $D$ has zero diagonal, is nonnegative, and is symmetric by the corresponding properties of $d$.\\
Fix $a,b,c\in\operatorname{CS}(X)$. Then $(\forall n\in\N)\,d(a_n,c_n)\leq d(a_n,b_n)+d(b_n,c_n)$ by triangle inequality.\\
Hence $D(a,c)\leq D(a,b)+D(b,c)$ by Proposition \ref{prop:convergent-sequence-basic-consequences}(iii) and Theorem \ref{thm:complex-sequence-limit-algebra}(i). Thus $D$ is a pseudometric.
\item \hlt{Quotient and embedding}: Define $a\sim b\Longleftrightarrow D(a,b)=0$ and set $\widehat{X}\defeq\operatorname{CS}(X)/{\sim}$. Put $\widehat{d}([a],[b])\defeq D(a,b)$.\\
Then $\sim$ is an equivalence relation and $\widehat{d}$ is a well-defined metric by Proposition \ref{prop:pseudometric-metric-quotient}.\\
For $x\in X$, let $\iota(x)$ be the class of the constant sequence with value $x$.\\
The identity $\widehat{d}(\iota(x),\iota(y))=d(x,y)$ makes $\iota:X\to\widehat{X}$ an isometry by Definition \ref{def:isometry-and-isometric-metric-spaces}(i).
\item \hlt{Density}: Fix $[a]\in\widehat{X}$ and $\varepsilon>0$. Choose $N_0\in\N$ with $(\forall m,n\geq N_0)\,d(a_m,a_n)<\frac{\varepsilon}{2}$.\\
For $N\geq N_0$, $\widehat{d}(\iota(a_N),[a])\leq\frac{\varepsilon}{2}<\varepsilon$ by Proposition \ref{prop:convergent-sequence-basic-consequences}(iii). Hence $\iota(a_N)\to[a]$, and $\overline{\iota(X)}^{\,\widehat{X}}=\widehat{X}$ by Proposition \ref{prop:sequential-closure-criterion}.
\item \hlt{Completeness}: Let $(\xi_k)$ be Cauchy in $\widehat{X}$. By (c), choose $x_k\in X$ with $\widehat{d}(\iota(x_k),\xi_k)<\frac{1}{k}$.\\
The bound $d(x_k,x_l)\leq\frac{1}{k}+\widehat{d}(\xi_k,\xi_l)+\frac{1}{l}$ follows by the triangle inequality.\\
Hence $(x_k)$ is Cauchy by Definition \ref{def:cauchy-sequence} and Theorem \ref{thm:special-sequence-limits}(i).\\
Set $a\defeq(x_k)\in\operatorname{CS}(X)$. Then $\widehat{d}(\xi_k,[a])<\frac{1}{k}+\lim\limits_{j\to\infty}d(x_k,x_j)$ by the triangle inequality.\\
Given $\varepsilon>0$, choose $K\in\N$ with $\frac{1}{K}<\frac{\varepsilon}{2}$ and $(\forall k,j\geq K)\,d(x_k,x_j)<\frac{\varepsilon}{2}$.\\
For $k\geq K$, the final limit is at most $\frac{\varepsilon}{2}$ by Proposition \ref{prop:convergent-sequence-basic-consequences}(iii). Hence $\widehat{d}(\xi_k,[a])<\varepsilon$.\\
Thus $\xi_k\to[a]$, and $(\widehat{X},\widehat{d},\iota)$ is a completion of $(X,d)$ by Definitions \ref{def:complete-metric-space} and \ref{def:metric-completion}.
\end{enumerate}
\item \hlt{Uniqueness}:
\begin{enumerate}[label=(\alph*)]
\setlength{\itemsep}{0pt}
\item \hlt{Dense-copy isometry}: Let $(Y,d_Y,\tau)$ and $(Z,d_Z,\sigma)$ be completions of $(X,d)$.\\
If $\tau x=\tau x'$, then $d(x,x')=d_Y(\tau x,\tau x')=0$, so $x=x'$; similarly, $\sigma$ is injective.\\
Define $\rho_0:\tau(X)\to\sigma(X)$ by $\rho_0(\tau x)\defeq\sigma x$. Iinjectivity of $\tau$ makes $\rho_0$ well defined, injectivity of $\sigma$ makes $\rho_0$ injective.\\
The identity $\rho_0(\tau x)=\sigma x$ implies $\rho_0(\tau(X))=\sigma(X)$, so $\rho_0$ is surjective.\\
For $x,x'\in X$, $d_Z(\rho_0(\tau x),\rho_0(\tau x'))=d(x,x')=d_Y(\tau x,\tau x')$. Hence $\rho_0$ is a bijective isometry by Definition \ref{def:isometry-and-isometric-metric-spaces}(i).
\item \hlt{Extension}: Fix $y\in Y$. Choose $(y_n)\subseteq\tau(X)$ with $y_n\to y$ by Proposition \ref{prop:sequential-closure-criterion}.\\
The sequence $(\rho_0(y_n))$ is Cauchy in $Z$ by the isometry property, so it converges by completeness of $Z$.\\
Define $\Phi(y)\defeq\lim\limits_{n\to\infty}\rho_0(y_n)$. Suppose $(z_n)\subseteq\tau(X)$ also satisfies $z_n\to y$.\\
The sequence $(\rho_0(z_n))$ is Cauchy by the isometry property and converges by completeness of $Z$.\\
Set $p\defeq\lim\limits_{n\to\infty}\rho_0(y_n)$ and $q\defeq\lim\limits_{n\to\infty}\rho_0(z_n)$. Then $d_Z(\rho_0(y_n),\rho_0(z_n))=d_Y(y_n,z_n)\leq d_Y(y_n,y)+d_Y(y,z_n)\to0$.\\
By triangle inequality, $d_Z(p,q)\leq d_Z(p,\rho_0(y_n))+d_Z(\rho_0(y_n),\rho_0(z_n))+d_Z(\rho_0(z_n),q)\to0$. Thus $p=q$, so $\Phi$ is well defined.
\item \hlt{Distance preservation}: For $y,y'\in Y$, choose $(y_n),(y_n')\subseteq\tau(X)$ converging to $y,y'$, respectively.\\
Set $u_n\defeq\rho_0(y_n)$, $v_n\defeq\rho_0(y_n')$. Inequality $\abs{d_Y(y_n,y_n')-d_Y(y,y')}\leq d_Y(y_n,y)+d_Y(y_n',y')$ implies $d_Y(y_n,y_n')\to d_Y(y,y')$.\\
$\abs{d_Z(u_n,v_n)-d_Z(\Phi(y),\Phi(y'))}\leq d_Z(u_n,\Phi(y))+d_Z(v_n,\Phi(y'))$ by triangle inequality. Hence $d_Z(u_n,v_n)\to d_Z(\Phi(y),\Phi(y'))$.\\
$(\forall n\in\N)\,d_Z(u_n,v_n)=d_Y(y_n,y_n')$. Hence $d_Z(\Phi(y),\Phi(y'))=d_Y(y,y')$ by Proposition \ref{prop:convergent-sequence-basic-consequences}(i).\\
Thus $\Phi$ is an isometry by Definition \ref{def:isometry-and-isometric-metric-spaces}(i).\\
For $x\in X$, the constant sequence $y_n=\tau x$ in the extension satisfies $\Phi(\tau x)=\sigma x$. Hence $\Phi\circ\tau=\sigma$.
\item \hlt{Surjectivity}: If $(\Phi(y_n))$ is Cauchy in $\Phi(Y)$, then $(y_n)$ is Cauchy in $Y$ and converges to some $y\in Y$.\\
The identity $d_Z(\Phi(y_n),\Phi(y))=d_Y(y_n,y)\to0$ implies $\Phi(y_n)\to\Phi(y)$, so $\Phi(Y)$ is complete.\\
Hence $\Phi(Y)$ is closed in $Z$ by Proposition \ref{prop:basic-completeness-consequences}(i).\\
The inclusions $\sigma(X)\subseteq\Phi(Y)$ and $Z=\overline{\sigma(X)}^{\,Z}$ imply $Z\subseteq\Phi(Y)$ by Theorem \ref{thm:closure-smallest-closed-containing-set}(iii).\\
Thus $\Phi$ is surjective; since every isometry is injective, it is bijective.
\item \hlt{Uniqueness}: Let $\Psi:Y\to Z$ be another isometry satisfying $\Psi\circ\tau=\sigma$, and fix $y\in Y$.\\
Choose $(y_n)\subseteq\tau(X)$ with $y_n\to y$ by Proposition \ref{prop:sequential-closure-criterion}.\\
Then $\Phi(y_n)=\Psi(y_n)$ and $d_Z(\Phi(y),\Psi(y))\leq d_Y(y,y_n)+d_Y(y_n,y)\to0$, so $\Phi(y)=\Psi(y)$. Hence $\Phi=\Psi$.
\end{enumerate}
\end{enumerate}
\hlt{Structural Commentary}: The quotient installs formal limits; density extends the copy of $X$ uniquely between completions.\\
\hlt{Key Insight}: Cauchy classes provide existence; dense agreement forces uniqueness.
\end{proof}
\end{theorem}

Figure \ref{fig:metric-completion-uniqueness} summarizes the unique extension between two completions.
\begin{illustrationfigure}[H]
  \centering
  % Source guide: Roman, Advanced Linear Algebra, Ch. 12, Figure 12.4.
% Redrawn schematic for uniqueness of metric completion.
\begin{tikzpicture}[
  >=Stealth,
  every node/.style={font=\small},
  source/.style={draw=orange!80!black, fill=orange!10, rounded corners=2pt,
    minimum width=1.35cm, minimum height=0.72cm, align=center, inner sep=5pt},
  completion/.style={draw=gray!60, fill=gray!7, rounded corners=2pt,
    minimum width=2.35cm, minimum height=0.76cm, align=center, inner sep=5pt},
  arrow/.style={->, line width=0.65pt, gray!75},
  iso/.style={->, line width=0.8pt, blue!70!black},
  note/.style={font=\scriptsize, fill=white, inner sep=1.2pt}
]
  \node[source] (X) at (0,0) {$X$};

  \node[completion] (Y) at (3.75,0.9) {$\tau(X)\subset Y$};
  \node[completion] (Z) at (3.75,-0.9) {$\sigma(X)\subset Z$};

  \draw[arrow] (X.east) -- node[note, above] {$\tau$} (Y.west);
  \draw[arrow] (X.east) -- node[note, below] {$\sigma$} (Z.west);
  \draw[iso] (Y.south) -- node[note, right] {$\overline{\rho}$} (Z.north);
\end{tikzpicture}

  \caption{The dense-copy isometry $\sigma\tau^{-1}$ extends uniquely to $\Phi:Y\cong Z$; adapted from \citet[Ch.~12, Fig.~12.4]{roman2008advanced}.}
  \label{fig:metric-completion-uniqueness}
\end{illustrationfigure}
